\subsection{Задание}

\begin{enumerate}
    \item Напишите программу для визуализации линий уровня целевой функции и ломаной последовательных приближений. Приближения итеративно порождаются с помощью методов оптимизации из ЛР4 до выполнения критерия останова. 
    
    Модифицируйте алгоритм отрисовки контуров так, чтобы линии уровня целевой функции проходили \textit{строго через вершины} ломаной приближений (соответственно, количество линий уровня должно быть равно количеству итераций/приближений). Постройте такую визуализацию для каждого метода оптимизации из ЛР4.

    \item Найдите точку минимума функции своего варианта численно \textbf{методом Ньютона}. В качестве критерия останова используйте порог на модуль градиента функции: $\|\nabla z\| \le 0.0001$.
    
    Выполните решение также \textit{аналитически} — найдите градиент, приравняйте его к нулю, определите все стационарные точки и с помощью матрицы Гессе укажите тип экстремума для каждой из них. Представьте в отчете подробную детализацию вычислений, а также реализуйте эти расчёты программно.

    \item Постройте \textbf{объектную модель} задачи минимизации функции двух переменных для трех методов из ЛР4: покоординатного спуска, градиентного спуска и метода наискорейшего спуска. 
    
    Выполните постановку задачи, выделите ключевые сущности (классы), опишите их атрибуты и методы. Спроектируйте иерархию классов (используя абстрактные классы и/или интерфейсы), которая объединяет логику всех трех методов. Воспользуйтесь объектными отношениями (наследование, реализация, композиция, агрегация или зависимость) по своему выбору. Обновите программную реализацию на языке Python с учетом разработанной объектной модели.
\end{enumerate}

\textbf{Исходные данные для варианта №3:}
\[
    z(x,y) = x^3 - 9x^2 + 24x + y^3 - 14y^2 + 64y - 116, \quad M_0 = (3.5, \, 5.5)
\]